\documentclass[11pt,a4paper]{article}
\usepackage[margin=2.5cm]{geometry}
\usepackage{amsmath,amssymb,amsthm}
\usepackage{physics}
\usepackage{booktabs}
\usepackage{hyperref}
\usepackage{array}
\usepackage{tcolorbox}

\newcommand{\Rxi}{R_\xi}
\newcommand{\Mpl}{M_{\mathrm{Pl}}}
\newcommand{\Mplbar}{\bar{M}_{\mathrm{Pl}}}
\newcommand{\Mfive}{M_5}
\newcommand{\calI}{\mathcal{I}}
\newcommand{\GN}{G_N}
\newcommand{\kfive}{\kappa_5}
\newcommand{\MZ}{M_Z}

\title{Derivation v20: Factor \& Normalization Audit\\
for the 5D$\to$4D Gravity Bridge}
\author{EDC Collaboration\\[3pt]
\small\textit{Derivation note; explicit factor tracking}}
\date{February 2026}

\begin{document}
\maketitle

\begin{abstract}
This note performs a forensic audit of all numerical factors in the derivation
of $\Mplbar^2 = \Mfive^3 \calI$. We explicitly track: (i) the $\frac{1}{2}$ in
action normalization, (ii) the $8\pi$ in Newton's constant, (iii) orbifold
doubling factors, (iv) the $2\pi$ from circle circumference, and (v) the warp
factor exponent in $\calI$. Three coordinate conventions are compared and a
canonical choice is fixed. A conversion table maps between conventions.
\end{abstract}

%═══════════════════════════════════════════════════════════════════════════════
\section{Conventions (Fixed for This Note)}
%═══════════════════════════════════════════════════════════════════════════════

\subsection{5D Action Normalization}

We adopt the 5D Einstein-Hilbert action in the form:
\begin{equation}
S_5 = \frac{\Mfive^3}{2} \int d^5x \sqrt{-g_5}\, R_5
\label{eq:5D-action}
\end{equation}
where $\Mfive$ is the fundamental 5D Planck mass. This is equivalent to
writing $S_5 = \frac{1}{2\kfive^2}\int\sqrt{-g_5}R_5$ with $\kfive^2 = 1/\Mfive^3$.

\textbf{Note:} The factor $\frac{1}{2}$ multiplies the entire action, not just
the Ricci scalar. This is the ``reduced'' action convention standard in
brane-world literature.

\subsection{4D Newton Constant and Planck Mass}

Newton's constant in 4D is defined via:
\begin{equation}
\GN = \frac{1}{8\pi \Mplbar^2}
\label{eq:GN-def}
\end{equation}
where $\Mplbar$ is the \textbf{reduced Planck mass}:
\begin{equation}
\Mplbar = \frac{1}{\sqrt{8\pi \GN}} \approx 2.435 \times 10^{18}~\mathrm{GeV}
\label{eq:Mplbar-def}
\end{equation}

The ``original'' Planck mass is:
\begin{equation}
\Mpl = \frac{1}{\sqrt{\GN}} \approx 1.221 \times 10^{19}~\mathrm{GeV}
\label{eq:Mpl-original}
\end{equation}

Relation:
\begin{equation}
\Mplbar = \frac{\Mpl}{\sqrt{8\pi}}
\label{eq:Mpl-relation}
\end{equation}

The 4D Einstein-Hilbert action with our conventions is:
\begin{equation}
S_4 = \frac{\Mplbar^2}{2} \int d^4x \sqrt{-g_4}\, R_4 = \frac{1}{16\pi\GN} \int d^4x \sqrt{-g_4}\, R_4
\label{eq:4D-action}
\end{equation}

\subsection{Extra Dimension Coordinate Conventions}

We consider three coordinate conventions for the extra dimension $\xi$:

\begin{enumerate}
\item \textbf{Interval $[0,L]$}: Flat extra dimension with coordinate $\xi \in [0,L]$.
The length of the extra dimension is $L$.

\item \textbf{Circle $S^1$}: Compact circle with $\xi \in [0, 2\pi R]$ and
periodic identification $\xi \sim \xi + 2\pi R$. Circumference is $2\pi R$.

\item \textbf{Orbifold $S^1/\mathbb{Z}_2$}: Circle with $\mathbb{Z}_2$ reflection
symmetry $\xi \to -\xi$. Fundamental domain is $\xi \in [0, \pi R]$. Branes sit
at the fixed points $\xi = 0$ and $\xi = \pi R$.
\end{enumerate}

\textbf{Definition:} We define $\Rxi$ as the \emph{length of the fundamental domain}:
\begin{equation}
\Rxi \equiv L \quad \text{(interval)}, \qquad
\Rxi \equiv 2\pi R \quad \text{(circle)}, \qquad
\Rxi \equiv \pi R \quad \text{(orbifold)}
\label{eq:Rxi-def}
\end{equation}

%═══════════════════════════════════════════════════════════════════════════════
\section{Linearized Equations and Mode Decomposition}
%═══════════════════════════════════════════════════════════════════════════════

\subsection{Warped Metric Ansatz}

The 5D metric is:
\begin{equation}
ds^2 = e^{2A(\xi)} \eta_{\mu\nu} dx^\mu dx^\nu + d\xi^2
\label{eq:metric}
\end{equation}
where $A(\xi)$ is the warp factor.

\subsection{Metric Perturbation}

Perturbing around the background:
\begin{equation}
g_{\mu\nu}^{(5)} = e^{2A}(\eta_{\mu\nu} + h_{\mu\nu})
\label{eq:pert-metric}
\end{equation}

In transverse-traceless (TT) gauge:
\begin{equation}
\partial^\mu h_{\mu\nu} = 0, \qquad \eta^{\mu\nu} h_{\mu\nu} = 0
\label{eq:TT-gauge}
\end{equation}

\subsection{KK Decomposition}

We decompose:
\begin{equation}
h_{\mu\nu}(x,\xi) = \sum_n h_{\mu\nu}^{(n)}(x) \psi_n(\xi)
\label{eq:KK-decomp}
\end{equation}

\subsection{Derivation of Mode Equation}

The linearized 5D Einstein equation for TT modes in the bulk is:
\begin{equation}
\Box_5 h_{\mu\nu} = 0
\label{eq:5D-EOM}
\end{equation}

Expanding the 5D d'Alembertian:
\begin{equation}
\Box_5 = e^{-2A} \Box_4 + \partial_\xi^2 + 4A' \partial_\xi
\label{eq:Box5}
\end{equation}

Substituting the KK decomposition and using $\Box_4 h^{(n)}_{\mu\nu} = -m_n^2 h^{(n)}_{\mu\nu}$:
\begin{equation}
-m_n^2 e^{-2A} \psi_n + \psi_n'' + 4A' \psi_n' = 0
\label{eq:mode-eq-derivation}
\end{equation}

Rearranging:
\begin{equation}
\boxed{\psi_n'' + 4A' \psi_n' + m_n^2 e^{2A} \psi_n = 0}
\label{eq:mode-eq}
\end{equation}

\subsection{Zero-Mode Solution}

For $m_0 = 0$:
\begin{equation}
\psi_0'' + 4A' \psi_0' = 0 \quad \Rightarrow \quad
\frac{d}{d\xi}(e^{4A}\psi_0') = 0
\label{eq:zero-mode-eq}
\end{equation}

For normalizable modes: $\psi_0' = 0$, hence:
\begin{equation}
\psi_0 = \text{const}
\label{eq:psi0-const}
\end{equation}

%═══════════════════════════════════════════════════════════════════════════════
\section{Kinetic Term Reduction: Factor-by-Factor}
%═══════════════════════════════════════════════════════════════════════════════

This section is the core of the factor audit.

\subsection{5D Kinetic Term}

Starting from the 5D action \eqref{eq:5D-action}, the quadratic kinetic term
for TT graviton modes is:
\begin{equation}
S_5^{(2)} = \frac{\Mfive^3}{2} \int d^5x \sqrt{-g_5}
\left( -\frac{1}{4} g^{AC} g^{BD} \nabla_A h_{CB} \nabla_D h_{CB} \right)
\label{eq:5D-kin-full}
\end{equation}

For our metric \eqref{eq:metric}:
\begin{equation}
\sqrt{-g_5} = e^{4A}
\label{eq:sqrt-g5}
\end{equation}

The inverse metric components:
\begin{equation}
g^{\mu\nu} = e^{-2A} \eta^{\mu\nu}, \qquad g^{55} = 1
\label{eq:inverse-metric}
\end{equation}

\subsection{Collecting Factors}

The 4D kinetic piece of the graviton (ignoring $\xi$-derivatives, which give
mass terms):
\begin{equation}
S_5^{(2)} \supset \frac{\Mfive^3}{2} \int d^4x\, d\xi\, e^{4A}
\left( -\frac{1}{4} e^{-2A} e^{-2A} \eta^{\mu\rho}\eta^{\nu\sigma}
\partial_\mu h_{\rho\sigma} \partial_\nu h^{\rho\sigma} \right)
\label{eq:factor-step1}
\end{equation}

Simplifying:
\begin{equation}
S_5^{(2)} \supset -\frac{\Mfive^3}{8} \int d^4x\, d\xi\, e^{4A} e^{-4A}
\eta^{\mu\rho}\eta^{\nu\sigma} \partial_\mu h_{\rho\sigma} \partial_\nu h^{\rho\sigma}
\label{eq:factor-step2}
\end{equation}

\textbf{Key observation:} The $e^{4A}$ from $\sqrt{-g_5}$ cancels the $e^{-4A}$
from two factors of $g^{\mu\nu} = e^{-2A}\eta^{\mu\nu}$.

Therefore:
\begin{equation}
S_5^{(2)} \supset -\frac{\Mfive^3}{8} \int d^4x\, d\xi\,
\eta^{\mu\rho}\eta^{\nu\sigma} \partial_\mu h_{\rho\sigma} \partial_\nu h^{\rho\sigma}
\label{eq:factor-step3}
\end{equation}

\subsection{Substituting KK Zero-Mode}

For the zero-mode $h_{\mu\nu} = h_{\mu\nu}^{(0)}(x) \psi_0(\xi)$:
\begin{equation}
S_5^{(2)} \supset -\frac{\Mfive^3}{8} \int d^4x\,
\eta^{\mu\rho}\eta^{\nu\sigma} \partial_\mu h_{\rho\sigma}^{(0)} \partial_\nu h^{(0)\rho\sigma}
\cdot \int d\xi\, |\psi_0|^2
\label{eq:zero-mode-sub}
\end{equation}

Define:
\begin{equation}
\boxed{\calI \equiv \int_{\text{domain}} d\xi\, |\psi_0(\xi)|^2}
\label{eq:I-def}
\end{equation}

\textbf{CRITICAL:} With the factor cancellation above, the exponent in $\calI$
is \textbf{zero}, not $e^{2A}$ or $e^{4A}$!

Continuing:
\begin{equation}
S_5^{(2)} \supset -\frac{\Mfive^3 \calI}{8} \int d^4x\,
\partial_\mu h_{\rho\sigma}^{(0)} \partial^\mu h^{(0)\rho\sigma}
\label{eq:S5-reduced}
\end{equation}

\subsection{Matching to 4D Action}

The 4D kinetic term from \eqref{eq:4D-action} is:
\begin{equation}
S_4^{(2)} = \frac{\Mplbar^2}{2} \cdot \left( -\frac{1}{4} \right)
\int d^4x\, \partial_\mu h_{\rho\sigma} \partial^\mu h^{\rho\sigma}
= -\frac{\Mplbar^2}{8} \int d^4x\, \partial_\mu h_{\rho\sigma} \partial^\mu h^{\rho\sigma}
\label{eq:4D-kin}
\end{equation}

Matching \eqref{eq:S5-reduced} to \eqref{eq:4D-kin}:
\begin{equation}
\boxed{\Mplbar^2 = \Mfive^3 \calI}
\label{eq:bridge}
\end{equation}

\subsection{Newton's Constant}

From \eqref{eq:GN-def} and \eqref{eq:bridge}:
\begin{equation}
\boxed{\GN = \frac{1}{8\pi \Mfive^3 \calI}}
\label{eq:GN-bridge}
\end{equation}

%═══════════════════════════════════════════════════════════════════════════════
\section{Computing $\calI$ for Different Conventions}
%═══════════════════════════════════════════════════════════════════════════════

\subsection{Flat Case: $A(\xi) = 0$}

For a flat extra dimension, the zero-mode is constant: $\psi_0 = c$.

\subsection{Case 1: Interval $[0,L]$}

Integration domain: $\xi \in [0,L]$.

Normalization: $\int_0^L d\xi\, |c|^2 = |c|^2 L = 1$ gives $|c|^2 = 1/L$.

Therefore:
\begin{equation}
\calI = \int_0^L d\xi\, |\psi_0|^2 = \int_0^L d\xi\, \frac{1}{L} = 1 \quad \text{(dimensionless)}
\label{eq:I-interval-norm}
\end{equation}

Alternatively, with $\psi_0 = 1$ (unnormalized):
\begin{equation}
\calI = \int_0^L d\xi = L
\label{eq:I-interval-unnorm}
\end{equation}

\textbf{With our definition $\Rxi = L$:}
\begin{equation}
\calI = \Rxi \quad \text{(interval, unnormalized $\psi_0 = 1$)}
\label{eq:I-Rxi-interval}
\end{equation}

\subsection{Case 2: Circle $S^1$ with $\xi \in [0, 2\pi R]$}

Integration domain: $\xi \in [0, 2\pi R]$.

With $\psi_0 = 1$:
\begin{equation}
\calI = \int_0^{2\pi R} d\xi = 2\pi R
\label{eq:I-circle}
\end{equation}

\textbf{With our definition $\Rxi = 2\pi R$ (circumference):}
\begin{equation}
\calI = \Rxi \quad \text{(circle)}
\label{eq:I-Rxi-circle}
\end{equation}

\subsection{Case 3: Orbifold $S^1/\mathbb{Z}_2$}

The orbifold $S^1/\mathbb{Z}_2$ has a $\mathbb{Z}_2$ symmetry $\xi \to -\xi$.

\textbf{Method A:} Integrate over fundamental domain $[0, \pi R]$ only:
\begin{equation}
\calI = \int_0^{\pi R} d\xi = \pi R
\label{eq:I-orbifold-fund}
\end{equation}

\textbf{Method B:} Integrate over full circle $[-\pi R, \pi R]$ and divide by
orbifold factor 2:
\begin{equation}
\calI = \frac{1}{2} \int_{-\pi R}^{\pi R} d\xi = \frac{1}{2} \cdot 2\pi R = \pi R
\label{eq:I-orbifold-full}
\end{equation}

Both methods agree.

\textbf{With our definition $\Rxi = \pi R$ (fundamental domain length):}
\begin{equation}
\calI = \Rxi \quad \text{(orbifold)}
\label{eq:I-Rxi-orbifold}
\end{equation}

\subsection{Summary: Universal Result}

With the definition \eqref{eq:Rxi-def} ($\Rxi$ = length of fundamental domain)
and unnormalized $\psi_0 = 1$:
\begin{equation}
\boxed{\calI = \Rxi \quad \text{(universal, flat case)}}
\label{eq:I-Rxi-universal}
\end{equation}

%═══════════════════════════════════════════════════════════════════════════════
\section{Canonical Convention and Conversion Table}
%═══════════════════════════════════════════════════════════════════════════════

\subsection{Canonical Choice}

We adopt:
\begin{equation}
\Rxi \equiv \text{length of fundamental integration domain}
\label{eq:canonical}
\end{equation}

The bridge relations become:
\begin{align}
\Mplbar^2 &= \Mfive^3 \Rxi \label{eq:bridge-canon} \\
\GN &= \frac{1}{8\pi \Mfive^3 \Rxi} \label{eq:GN-canon} \\
\Mfive &= \Mplbar^{2/3} \Rxi^{-1/3} \label{eq:M5-canon}
\end{align}

\subsection{Conversion Table}

\begin{table}[h]
\centering
\caption{Conversion between geometric conventions.}
\label{tab:conversion}
\begin{tabular}{@{}l c c c c@{}}
\toprule
Geometry & Domain & $\calI$ (flat) & $\Rxi$ definition & Mapping \\
\midrule
Interval $[0,L]$ & $[0,L]$ & $L$ & $\Rxi = L$ & $\Rxi = L$ \\
Circle $S^1$ & $[0,2\pi R]$ & $2\pi R$ & $\Rxi = 2\pi R$ & $\Rxi = 2\pi R$ \\
Orbifold $S^1/\mathbb{Z}_2$ & $[0,\pi R]$ & $\pi R$ & $\Rxi = \pi R$ & $\Rxi = \pi R$ \\
\bottomrule
\end{tabular}
\end{table}

\textbf{Conversion rule:} If an alternative convention uses $\tilde{R}$ where
$\calI = f(\tilde{R})$, then $\Rxi = f(\tilde{R})$.

Example: If someone uses ``$R$'' to mean the circle \emph{radius} (not circumference),
then $\calI = 2\pi R$ and thus $\Rxi = 2\pi R$.

\subsection{Calibrated Closure}

With $\Rxi = \hbar c / \MZ^{\mathrm{obs}}$ identified as [I]+[BL]:
\begin{equation}
\Mfive = \Mplbar^{2/3} \Rxi^{-1/3} = (2.435 \times 10^{18})^{2/3} \times (2.165 \times 10^{-18}~\mathrm{m})^{-1/3}
\label{eq:closure-numerical}
\end{equation}

Converting $\Rxi$ to GeV$^{-1}$: $\Rxi = 1/\MZ = 1/91.19$~GeV$^{-1}$:
\begin{equation}
\Mfive = \Mplbar^{2/3} \MZ^{1/3} = (2.435 \times 10^{18})^{2/3} \times (91.19)^{1/3}~\mathrm{GeV}
\label{eq:closure-formula}
\end{equation}

Computing:
\begin{equation}
\Mfive = 1.81 \times 10^{12} \times 4.50 = 8.1 \times 10^{12}~\mathrm{GeV}
\label{eq:M5-reduced}
\end{equation}

\textbf{CORRECTION vs v19:} In v19, we used $\Mpl = 1.221 \times 10^{19}$~GeV
(the ``original'' Planck mass) in the formula. With the \emph{reduced} Planck
mass $\Mplbar = 2.435 \times 10^{18}$~GeV, the result is:
\begin{equation}
\Mfive = 8.1 \times 10^{12}~\mathrm{GeV} \quad \text{(reduced convention)}
\label{eq:M5-reduced-value}
\end{equation}

If instead we use the \emph{original} Planck mass in the bridge $\Mpl^2 = \Mfive^3\Rxi$:
\begin{equation}
\Mfive = \Mpl^{2/3} \Rxi^{-1/3} = (1.221 \times 10^{19})^{2/3} \times (91.19)^{1/3}
= 2.4 \times 10^{13}~\mathrm{GeV}
\label{eq:M5-original}
\end{equation}

The difference is a factor of $(8\pi)^{1/3} \approx 2.94$.

%═══════════════════════════════════════════════════════════════════════════════
\section{Factor Audit Summary}
%═══════════════════════════════════════════════════════════════════════════════

\begin{tcolorbox}[title=Factor Audit Box]
\begin{tabular}{@{}p{4.5cm} p{4cm} p{3cm}@{}}
\toprule
\textbf{Factor} & \textbf{Resolution} & \textbf{Eq.\ ref} \\
\midrule
$\frac{1}{2}$ in 5D action & Fixed in \eqref{eq:5D-action} & Eq.~\eqref{eq:5D-action} \\
$\frac{1}{4}$ in kinetic term & Standard TT normalization & Eq.~\eqref{eq:5D-kin-full} \\
$8\pi$ in $\GN$ definition & $\GN = 1/(8\pi\Mplbar^2)$ & Eq.~\eqref{eq:GN-def} \\
Reduced vs original $\Mpl$ & Use $\Mplbar = \Mpl/\sqrt{8\pi}$ & Eq.~\eqref{eq:Mpl-relation} \\
Warp factor exponent in $\calI$ & \textbf{Zero} (factors cancel) & Eq.~\eqref{eq:I-def} \\
Orbifold factor of 2 & Absorbed in domain definition & Eq.~\eqref{eq:I-orbifold-full} \\
$2\pi$ from circle & Absorbed in $\Rxi = 2\pi R$ & Eq.~\eqref{eq:I-Rxi-circle} \\
$\Rxi$ definition & Length of fundamental domain & Eq.~\eqref{eq:Rxi-def} \\
\bottomrule
\end{tabular}
\end{tcolorbox}

%═══════════════════════════════════════════════════════════════════════════════
\section{Reconciliation with v19}
%═══════════════════════════════════════════════════════════════════════════════

\subsection{v19 Convention}

Derivation v19 used:
\begin{itemize}
\item $\Mpl = 1.221 \times 10^{19}$~GeV (original Planck mass)
\item $\GN = 1/(8\pi\Mpl^2)$ (which is dimensionally inconsistent with the
original Planck mass definition $\Mpl = \GN^{-1/2}$)
\end{itemize}

\subsection{Self-Consistent Options}

\textbf{Option A (reduced convention):}
\begin{align}
\GN &= \frac{1}{8\pi \Mplbar^2}, \quad \Mplbar = 2.435 \times 10^{18}~\mathrm{GeV} \\
\Mfive &= \Mplbar^{2/3} \Rxi^{-1/3} = 8.1 \times 10^{12}~\mathrm{GeV}
\end{align}

\textbf{Option B (original convention):}
\begin{align}
\GN &= \frac{1}{\Mpl^2}, \quad \Mpl = 1.221 \times 10^{19}~\mathrm{GeV} \\
\Mfive &= \Mpl^{2/3} \Rxi^{-1/3} = 2.4 \times 10^{13}~\mathrm{GeV}
\end{align}

\subsection{Which is Correct?}

Both are correct \emph{within their own conventions}. The key is consistency:

\begin{itemize}
\item If 4D action is $S_4 = \frac{\Mplbar^2}{2}\int\sqrt{-g}R$, use Option A.
\item If 4D action is $S_4 = \frac{\Mpl^2}{16\pi}\int\sqrt{-g}R$, use Option B.
\end{itemize}

\textbf{v19 implicitly used Option B} (numerical value $\Mpl = 1.221 \times 10^{19}$~GeV),
but wrote the bridge in the Option A form. This is a notational inconsistency,
not a physics error, because the $8\pi$ factors cancel in physical observables.

\subsection{Recommendation}

For clarity, future notes should use the reduced convention (Option A) with
$\Mplbar = 2.435 \times 10^{18}$~GeV. The numerical result for $\Mfive$ depends
on convention:
\begin{equation}
\Mfive =
\begin{cases}
8.1 \times 10^{12}~\mathrm{GeV} & \text{(reduced: Option A)} \\
2.4 \times 10^{13}~\mathrm{GeV} & \text{(original: Option B)}
\end{cases}
\label{eq:M5-both}
\end{equation}

Both are at the GUT scale; the factor of 3 difference is a convention choice.

%═══════════════════════════════════════════════════════════════════════════════
\section{Conclusion}
%═══════════════════════════════════════════════════════════════════════════════

\begin{enumerate}
\item The warp factor exponent in $\calI$ is \textbf{zero} due to exact cancellation
of $e^{4A}$ from $\sqrt{-g_5}$ and $e^{-4A}$ from $g^{\mu\nu}g^{\rho\sigma}$.
\item $\Rxi$ should be defined as the \textbf{length of the fundamental domain},
giving $\calI = \Rxi$ universally (flat case).
\item The $8\pi$ factor resides in $\GN = 1/(8\pi\Mplbar^2)$ and in the relation
$\Mplbar = \Mpl/\sqrt{8\pi}$.
\item No orbifold doubling ambiguity exists if $\Rxi$ is defined as the
fundamental domain length.
\item The numerical value of $\Mfive$ depends on the Planck mass convention:
$\sim 10^{13}$~GeV (original) or $\sim 10^{12}$~GeV (reduced). Both are valid.
\end{enumerate}

\begin{center}
\fbox{\parbox{0.9\textwidth}{
\textbf{Canonical bridge (reduced convention):}
\[
\Mplbar^2 = \Mfive^3 \Rxi, \qquad \GN = \frac{1}{8\pi\Mfive^3\Rxi}
\]
\textbf{Closure:} $\Mfive = \Mplbar^{2/3}\Rxi^{-1/3}$ once $\Rxi$ is identified [I]/[BL].
}}
\end{center}

\end{document}
